\section{循环过程}
\label{sec:循环过程}
在历史上，热力学理论的建立与发展，是与热机以及热泵的使用和改进密切相关的，而这两者的实际原理虽然很复杂，但就其工作模式而言，可以抽象为本节将介绍的热力学循环过程。

\subsection{热机和热泵}
\label{subsec:热机和热泵}
我们先来看两个简单的事实，气体是具有能量的，即，气体的内能。而在另外一方面，气体的能量可以通过功和热与外界进行交换。在一个准静态过程中，气体的内能可能发生变化，但是如果我们使得内能在过程始末保持不变，也就是说使得气体经过一系列变化过程回到原来的状态，\empx{那么我们就有可能持续的进行功和热间的转化}，因为这两者是过程量，即便气体恢复到原来的状态功和热各自的净值也未必为零，但两者的和必然为零，根据$\delt{E}=W+Q=0$。

我们在上述讨论中可以看出，使气体在每一次准静态过程后恢复原有状态是一个要点，若气体的状态不能恢复，那么功和热间的转化就不能反复进行。这种特殊的过程称为循环过程。
\begin{OldBoxDefinition}[循环过程]
    如果一个热力学系统由某个状态出发，经过任意的一系列准静态变化过程，最终又回到原来的状态，那么，我们就将这样的过程称为\uwave{循环过程}，循环过程的一个重要特征是
    \begin{Equation}&[]
        \delt{E}=0
    \end{Equation}
    即循环过程前后，热力学系统的内能不变。
\end{OldBoxDefinition}

循环过程可以用P--V图上的一条闭合曲线描述，因此很明显，循环过程具有方向性\footnote{在先前热力学的内容中，我们将$W$的意义规定为“外界对系统做功”，这是因为在当时我们比较关心的是系统的能量是增加还是减少，但是，在循环过程中，我们更关心的是系统对外的效应，因此更多采用“系统对外做功”的表述，使用$W'$以避免混淆。}
\begin{OldFigure}[循环过程]
    \begin{OldFigureSub}[正循环]
        \inputfigure[scale=0.9]{Chapter12D_Figure02}
    \end{OldFigureSub}
    \hspace{1cm}
    \begin{OldFigureSub}[逆循环]
        \inputfigure[scale=0.9]{Chapter12D_Figure01}
    \end{OldFigureSub}
\end{OldFigure}

循环过程的曲线总是可以分为膨胀和压缩两部分，而我们知道P--V图上的曲线下面积代表功，以对外做功$W'$记，膨胀为正，压缩为负，但是由于膨胀和压缩过程未必在同一曲线上，因此两者相减的净值，即两者构成的闭合曲线间的面积，就是循环过程中对外做功$W'$的净值。
\begin{itemize}
    \item 如\xref{fig:正循环}所示的是\uwave{正循环过程}，该过程中，膨胀过程处于高位，压缩过程处于低位，这就表明，正循环过程中$W'$为正，系统将对外界做功，系统会吸收热量保证能量守恒。
    \item 如\xref{fig:逆循环}所示的是\uwave{逆循环过程}，该过程中，膨胀过程处于低位，压缩过程处于高位，这就表明，逆循环过程中$W'$为负，外界将对系统做功，系统会放出热量保证能量守恒。
\end{itemize}
从热力学第一定律上看，正循环即吸热对外做正功，逆循环即放热对外做负功，然而，实际情况并没有这么的简单，出于种种原因，循环过程中系统与外界的热量交换不可能或通常不会仅与单一热源进行。\empx{对于正循环而言这是不可能}，热力学第二定律会会告诉我们，尽管能量上守恒，但我们不能从单一热源吸收热量并完全转化为功。\empx{对于逆循环而言这是通常不会}，将功转化为热是很容易的事情（如电热水壶和电磁炉），但其实逆循环还可以实现更多的事情。

总而言之，\empx{通常的循环过程都是在两个温度不同的热源间进行的}，即存在三个部分，首先是实现循环过程的工作介质，例如一定量的理想气体，其次就是一个\uwave{高温热源}和\uwave{低温热源}，它们将在循环过程中分别与工作介质发生热量交换。正循环和逆循环的差异就在于热量交换的方向
\begin{OldFigure}[循环过程中的能量流动]
    \begin{OldFigureSub}[热机的工作模式]
        \inputfigure[scale=0.78]{Chapter12D_Figure03}
    \end{OldFigureSub}
    \begin{OldFigureSub}[热泵的工作模式]
        \inputfigure[scale=0.78]{Chapter12D_Figure04}
    \end{OldFigureSub}
\end{OldFigure}
\begin{OldBoxDefinition}[热机]
    \uwave{热机}是基于正循环工作，通过吸收热量实现对外做功的机器。

    热机从高温热源中吸收的热量，部分用于对外做功，部分放出至低温热源。
\end{OldBoxDefinition}
热机在现实中的具体例子就是\uwave{外燃机}（主要是\uwave{蒸汽机}）和\uwave{内燃机}，两者的主要差别在于热源是否在工质外部，蒸汽机的工质是蒸汽，其加热需要通过蒸汽机外部的燃料实现，内燃机的工质就是汽油或柴油一类的燃料，其自身就可以被点燃，两者各有优劣，在不同场合均有应用。

热机不可能完全的将从高温热源吸收的热量转化为对外的机械功，必然会在低温热源上以热量形式散失一部分的能量，因此衡量热机的一个重要指标就是其能多大比例的将热转化为功。
\begin{OldBoxDefinition}[热机效率]
    热机的效率是其热功转化的比例，即
    \begin{Equation}
        \eta=\frac{W'}{Q_1}=\frac{Q_1-Q_2}{Q_1}=1-\frac{Q_2}{Q_1}
    \end{Equation}
    其中$W'$表示热机对外界做的功。
\end{OldBoxDefinition}

\begin{OldBoxDefinition}[热泵]
    热泵是基于逆循环工作，通过外界做功实现热量逆向输运的机器。

    热泵向高温热源中放出的热量，部分来自外界做功，部分来自低温热源。
\end{OldBoxDefinition}
热泵的工作方式与热机刚好相反，在一定意义上
\begin{itemize}
    \item 热机是从温度差上的热量流动中截取能量，相当于水力发电机。
    \item 热泵是通过外加能量使热量逆向流动，相当于水泵将水从低处向高出运输。
\end{itemize}

热泵在现实中的具体例子就是空调或电冰箱。我们知道，热泵的效果是使得热量由低温热源流向高温热源，热机的效果是对外做功，因此不同于热机单一的工作方式，热泵具体取决于我们的温度调控目标是低温热源还是高温热源，会产生两种不同的工作方式，以空调为例的话
\begin{itemize}
    \item 对于冷空调，室内是低温热源，室外是高温热源，热量由室内“抽至”室外，实现降温。
    \item 对于热空调，室内是高温热源，室外是低温热源，热量由室外“抽至”室内，实现升温。
\end{itemize}
这就产生了一个很有意思的结果，冷空调降低室内温度时实际上还从室内获取了热量，但我们甚至还要为此付出额外的能量。然而，热空调则有所不同，热空调用于升高室内温度的能量其实主要来自室外，而我们施加的能量只是起到了一个泵的作用，当然这额外的小部分能量最终也转化为了室内的热，由此可见，\empx{热的产生总是容易的}。这也就解释了一个困惑，既然空调需要用电，那为什么我们不直接烧电阻丝来实现热空调的功能？这就是因为如果我们将这些电功转化为机械功，使其驱动热泵的工作，我们就可以从外部的低温环境中获得远多于机械功的热量，即，\empx{热泵的工作可以利用环境热量，产生比施予的外部能量更多的热量}，当然现代的空调在制热时为了快速升高温度，往往会带有“辅热”功能，即通过烧电阻丝迅速产生热量。

\begin{OldBoxDefinition}[热泵效能]
    热泵在制热时的效能是其放出的热量与投入的功的比例，即
    \begin{Equation}
        \varepsilon_{H}=\frac{Q1}{W'}=\frac{Q_1}{Q_1-Q_2}
    \end{Equation}
    热泵在制冷时的效能是其吸收的热量与投入的功的比例，即
    \begin{Equation}
        \varepsilon_{C}=\frac{Q2}{W'}=\frac{Q_2}{Q_1-Q_2}
    \end{Equation}
    其中$W$表示热泵由外界提供的功。
\end{OldBoxDefinition}
热泵的制热效能$\varepsilon_H$刚好是热机效率$\eta$的倒数，因为\empx{热泵在制热意义下即反向运作的热机}。

那么现在就产生了一些问题，我们已经知道如何从实际过程中的热和功去衡量热机和热泵的效率和效能，但是，如何从理论上去计算热机和热泵的效率和效能？如何确定效率和效能的理论上限（或者说效率和效能是否存在某个上限）？这就是我们将在下一小节中解决的问题。

\subsection{卡诺循环}
\label{subsec:卡诺循环}
在1824年法国青年工程师卡诺提出了一种理想模型，假设工作物质仅在高温和低温两个恒温热源间交换热量，在与热源交换热量时作等温变化，在两个热源间切换时作绝热变化，这种循环过程称为\uwave{卡诺循环}，取决于循环的正逆，该模型被称为\uwave{卡诺热机}或\uwave{卡诺热泵}。并且，卡诺还证明了卡诺循环具有最高的效率或效能。接下来我们以卡诺热机为例介绍卡诺循环的过程。
\begin{OldFigure}[卡诺循环]
    \hspace{0.2cm}
    \begin{OldFigureSub}[卡诺热机]
        \inputfigure[scale=0.7]{Chapter12D_Figure05}
    \end{OldFigureSub}
    \hspace{0.5cm}
    \begin{OldFigureSub}[卡诺热泵]
        \inputfigure[scale=0.7]{Chapter12D_Figure06}
    \end{OldFigureSub}
\end{OldFigure}
卡诺循环过程分为四步，包括两个等温过程和两个绝热过程
\begin{enumerate}
    \item 气缸与温度为$T_1$的高温热源接触，吸收热量$Q_1$，等温膨胀，对外做功。
    \item 气缸脱离高温热源，绝热膨胀，对外做功，温度随之下降至$T_2$。
    \item 气缸与温度为$T_2$的低温热源接触，放出热量$Q_2$，等温压缩，外界做功。
    \item 气缸脱离低温热源，绝热压缩，外界做功，温度随之升高至$T_1$。
\end{enumerate}
卡诺正循环在P--V图中表现为A$\to$B$\to$C$\to$D的循环过程，如\xref{fig:卡诺热机}所示，直观的说，卡诺循环实际就是通过绝热线在两条温度不同的等温线间反复切换的过程。卡诺循环中等温过程和绝热过程均有做功，但是，卡诺循环仅在等温过程交换热量，这就为计算效率带来了便利。

设高温热源吸热$Q_1$，设低温热源放热$Q_2$，根据\fancyref{ppt:等温过程}
\gdef\prefixeq{卡诺热机}
\begin{Gather}[8pt]
    Q_1=\nu RT_1\ln\frac{V_\text{B}}{V_\text{A}}\xlabelpeq{1}\\
    Q_2=\nu RT_2\ln\frac{V_\text{C}}{V_\text{D}}
    \xlabelpeq{2}
\end{Gather}
注意到\xrefpeq{2}中是$\ln V_\text{C}/V_{\text{D}}$而不是正常的$\ln V_\text{D}/V_{\text{C}}$，这是因为$Q_2$被设为放热而不是吸热。

根据\fancyref{def:热机效率}
\begin{Equation}&[3]
    \eta=1-\frac{Q_2}{Q_1}=
    1-\frac
    {\nu RT_2\ln(V_\text{C}/V_\text{D})}
    {\nu RT_1\ln(V_\text{B}/V_\text{A})}
\end{Equation}
应用\fancyref{fml:泊松公式}可以使得\xrefpeq{3}得到进一步的简化。

对于绝热膨胀过程B$\to$C有
\begin{Equation}&[4]
    T_1V_\text{B}^{\gamma-1}=T_2V_\text{C}^{\gamma-1}
\end{Equation}

对于绝热压缩过程D$\to$A有
\begin{Equation}&[5]
    T_2V_\text{D}^{\gamma-1}=T_1V_\text{A}^{\gamma-1}
\end{Equation}

比较\xrefpeq{4}和\xrefpeq{5}即有
\begin{Equation}&[6]
    \frac{V_\text{B}}{V_\text{A}}=
    \frac{V_\text{C}}{V_\text{D}}
\end{Equation}
这样，将\xrefpeq{6}代入\xrefpeq{3}中即得
\begin{OldBoxFormula}[卡诺热机的效率]
    卡诺热机的效率为
    \begin{Equation}
        \eta=\frac{T_1-T_2}{T_1}=1-\frac{T_2}{T_1}
    \end{Equation}
    卡诺热机的效率仅由高温和低温热源的温度决定，温度比越大，效率越高。
\end{OldBoxFormula}
这就是说，如果要提高热机的效率，我们就要使得高温热源和低温热源间的温度比尽可能大。

\empx{卡诺热机被证明是理论上效率最高的热机}，也就是说，在高温热源$T_1$和低温热源$T_2$之间工作的任何热机，即便在没有损耗的理想条件下，也不可能具有超过卡诺热机$\eta=1-(T_2/T_1)$的效率，这就是\uwave{卡诺定理}，它为热机的改进提供了理论指导，并给出了热机效率的理论上限。

卡诺逆循环与卡诺正循环是很相似的，只不过由原先A$\to$B$\to$C$\to$D变为了D$\to$C$\to$B$\to$A的相反方向，为了计算卡诺热泵的制热效能和制冷效能，我们只需要重复以上过程，不难得出
\begin{OldBoxFormula}[卡诺热泵的效能]
    卡诺热泵的制热效能是
    \begin{Equation}
        \varepsilon_H=\frac{T_1}{T_1-T_2}
    \end{Equation}
    卡诺热泵的制冷效能是
    \begin{Equation}
        \varepsilon_H=\frac{T_2}{T_1-T_2}
    \end{Equation}
    卡诺热泵的效能仅由高温和低温热源的温度决定，两者与温差的比越大，相应效能越高。
\end{OldBoxFormula}
这就是说，如果温差已经很大了，那么进一步增大温差就需要更多的功。